\chapter{Skinning and Animation across Two Systems}
\index{skinning}\index{animation}
\label{ch:skinning-animation}

CNA contains two skeletal-animation systems that deliberately do not share a runtime graph.
\cnaclass{SkinningData} plus \cnaclass{AnimationPlayer} animate an ordinary
\cnaclass{Model}; \cnaclass{SkinnedModelEXT} is an Avatar-oriented model, skeleton, clip, and
draw system of its own. Similar interpolation code appears in both. The duplication is
documented as intentional because forcing the public types together would couple two distinct
asset contracts.

\begin{figure}[tbp]
\centering
\begin{figurealt}{Three index spaces are connected by explicit maps. Flattened selected-scene nodes become ModelBone indices. The ordinal in the glTF skin joints array is the source joint index used by JOINTS zero. BuildSkeleton topologically reorders joints and records an old-to-new map, yielding palette indices. Packed joint values are rewritten to palette indices before SkinningData and the skinned effect consume them.}
\resizebox{0.96\textwidth}{!}{%
\begin{tikzpicture}[node distance=10mm and 14mm]
  \node[cna public] (scene) {selected-scene node\\index};
  \node[cna public, below=of scene] (joint) {source-joint ordinal\\\texttt{skin.joints}/\texttt{JOINTS\_0}};
  \node[cna internal, right=20mm of joint] (map) {topological reorder\\old $\rightarrow$ new map};
  \node[cna public, right=of map] (palette) {palette index};
  \node[cna internal, above=of palette] (bone) {\texttt{ModelBone} index\\from flattened scene};
  \node[cna internal, right=of palette] (skin) {bind + inverse bind +\\animated worlds};
  \node[cna public, right=of skin] (effect) {\texttt{SkinnedEffect} palette};
  \node[cna internal, below=of map] (rewrite) {rewrite packed joint bytes};
  \draw[cna flow] (scene) -- (bone);
  \draw[cna flow] (joint) -- (map);
  \draw[cna flow] (map) -- (palette);
  \draw[cna flow] (joint) -- (rewrite);
  \draw[cna flow] (map) -- (rewrite);
  \draw[cna flow] (rewrite) -- (palette);
  \draw[cna flow] (palette) -- (skin);
  \draw[cna flow] (skin) -- (effect);
  \draw[cna optional] (bone) -- node[cna label]{placement relation} (skin);
\end{tikzpicture}}
\end{figurealt}
\caption{Scene-node, source-joint, and palette indices are distinct. Convenient source ordering
can hide a missing conversion, so the importer records and tests the maps explicitly.}
\label{fig:skin-index-spaces}
\end{figure}



\section{The ordinary Model path}

A skinned \cnaclass{Model} carries \cnaclass{SkinningData} through \texttt{Model.Tag}. The
record contains bind-pose local matrices, inverse-bind matrices, parent indices, animation
clips, and a fourth matrix array named \texttt{SkeletonRootPrefix}. An
\cnaclass{AnimationPlayer} samples a clip into local transforms, composes worlds in topological
order, then builds the skin palette used by \cnaclass{SkinnedEffect} or
\cnaclass{SkinnedPbrEffect}.

The model itself does not do that last step. \texttt{Model::Draw} sets the ordinary
\cnaclass{IEffectMatrices} world matrix for each mesh, but never installs a bone palette. The
application or animation helper must call the skinned effect's
\texttt{SetBoneTransforms} before drawing. A model whose effect does not implement
\cnaclass{IEffectMatrices} is rejected by \texttt{Model::Draw}.

\section{The Avatar path}

\cnaclass{SkinnedModelEXT} is explicitly not built from \cnaclass{ModelBone},
\cnaclass{ModelMesh}, or \cnaclass{ModelMeshPart}. It owns its Avatar-specific geometry and
animation representation and is loaded through the legacy \texttt{.skinnedmodel.json} reader.
Its superficially similar \texttt{SampleTrack} logic is a local copy, not a hidden conversion
to \cnaclass{AnimationPlayer}.

Both players nevertheless enforce the same essential inputs. Bone arrays must agree in size,
and every non-root parent index must precede its child; a parent index greater than or equal to
the current index throws. Any negative parent value is treated as ``root'', not only $-1$, and an
animation track with an empty key list or an out-of-range bone index is silently ignored. Looping
uses floor-modulo on raw ticks, so negative or over-duration times wrap consistently. Runtime
translations and scales interpolate linearly, rotations spherically. Track search is linear in
the number of keyframes per bone.

The size agreement is not a complete hand-construction validator. Content readers reject a
negative or implausibly large bone count, but \cnaclass{AnimationPlayer}'s public constructor first
casts \texttt{BoneCount} to \texttt{size\_t} for its working-vector allocations and only then runs
the array-size check in \texttt{RecomputeTransforms}. A manually supplied negative count can
therefore request an enormous allocation (typically failing with a standard allocation/length
exception) rather than producing the later \cnaclass{ArgumentException}. Validate custom
\cnaclass{SkinningData} as non-negative before construction.

\section{Scene indices are not palette indices}

The glTF scene is flattened into \cnaclass{ModelBone} indices, while a skin's
\texttt{joints} array defines a separate palette. Those spaces need not have the same order and
the glTF joint array need not be parent-before-child. \texttt{BuildSkeleton} therefore
reorders joints breadth-first into topological order and records an old-to-new map. Packed
\texttt{JOINTS\_0} values are rewritten through that map.

This yields three indices worth naming explicitly:

\begin{description}
  \item[scene-node index] selects the model bone representing placement in the selected scene;
  \item[source-joint index] is the ordinal stored in glTF \texttt{JOINTS\_0}; and
  \item[palette index] selects the reordered bind, inverse-bind, and animated world arrays.
\end{description}

Treating any two as interchangeable works only on conveniently authored fixtures. Reversed
joint-order tests exist precisely to rule out that accident.

\section{The coordinate-space failure that produced D8}

A joint root may have ancestors that are not themselves listed in \texttt{skin.joints}. Older
logic walked parents only within the joint set. It therefore dropped the armature node's world
transform from the reconstructed bind hierarchy while retaining the file's inverse-bind matrix,
which already encoded that ancestry. A uniform armature scale became an inverse scale in the
final palette, collapsing the character toward the origin.

The corrected root term is

\[
  P_{root} = W_{ancestor} \, W_{meshNode}^{-1}.
\]

The first factor recovers scene ancestry above the joint set. The second cancels the transform
of the node that instances the skinned mesh, because the mesh is parented to the model's
synthetic root rather than to that node. The importer does not stop at glTF's optional
\texttt{skin.skeleton} hint, and can obtain ancestry for a joint outside the selected scene
through cgltf's world-transform calculation.

\subsection{Why the prefix must remain separate}

It is tempting to multiply $P_{root}$ into the root bind-pose local matrix. That works until an
animation channel replaces that local transform. At the first sampled frame the recovered
ancestry disappears again. Instead, CNA stores one prefix per root in
\texttt{SkeletonRootPrefix} and composes

\[
  W_{root}(t) = L_{root}(t) P_{root}.
\]

An empty prefix array means identity, preserving older CNJ skeleton sidecars. This separation
is both the D8 fix and the compatibility mechanism that lets the runtime animate roots without
undoing it.

The conformance fixture for the mesh-node term was chosen so the three possible results are
numerically distinct: missing cancellation, cancellation exactly once, and applying the
inverse twice. The test expects the middle case rather than merely checking that a vertex
``moved''.

\section{Clip extraction}

For skeletal clips, CNA forms the union of translation, rotation, and scale key times for each
animated bone, deduplicated at a tight tolerance. Missing channels at a time fall back to the
decomposed bind pose. Linear and step samplers are evaluated directly; cubic splines use the
glTF Hermite basis with the required $\Delta t$ tangent scaling, and cubic quaternion results
are normalized \emph{while the importer samples that union}. The resulting ordinary
\cnaclass{KeyframeEXT} records do not retain bone-channel interpolation modes or tangents;
\cnaclass{AnimationPlayer} later uses linear/Slerp interpolation between them. A cubic channel is
therefore exact at every union time but generally only a piecewise linear/Slerp approximation
between those stored times. Morph tracks below are different: they retain cubic tangents for lazy
Hermite evaluation.

Only channels whose target resolves through the skin's joint map enter a skeletal clip. In
addition, both the runtime and offline front ends call skeletal clip extraction only for a
group that has a skin. Rigid node TRS animation is thus dropped twice and without a warning.
This is the open D6 defect. The existence of correct interpolation code must not be reported as
general scene-animation support.

Morph-weight animation follows a different, skin-independent route. It locates a weights
channel on the node instancing the mesh, preserves cubic in/out tangents unbaked, and evaluates
Hermite weights lazily during playback. If one mesh is instanced by several nodes, the first
matching node wins. Morph position/normal deltas are then CPU-blended into vertex bytes and the
buffer is uploaded again; this is not GPU morphing.

\section{Scaling and capacity}

The converter's \texttt{unitScale} applies to every translation-bearing quantity, including
animation translation values and their cubic tangents. It never applies to rotation or scale
channels. Runtime glTF import fixes it at~1.0, while the CLI can choose another unit basis.

Both stock skinned effects accept at most 72 palette matrices. Each packed vertex has four
8-bit joint indices and four weights. A legal glTF skin larger than those runtime limits must
be split or otherwise reduced during content processing; successful container parsing does not
promise it can be drawn in one part.

\section{Practical update order}

A frame using the ordinary path should perform these operations in order:

\begin{enumerate}
  \item choose and sample the active clip into local bone transforms;
  \item compose absolute skeleton transforms, including root prefixes;
  \item form skin matrices from worlds and inverse binds;
  \item install the palette on every skinned effect used by the model;
  \item evaluate and upload morph weights if they changed; and
  \item call \texttt{Model::Draw} on the graphics thread.
\end{enumerate}

Do not concurrently draw models through \texttt{Model::Draw}: its bone-matrix scratch vector is
process-wide and unsynchronized. This restriction is independent of which animation player
produced the palette.
